\usepackage{xspace}
\usepackage{xcolor}
\usepackage{soul}
\usepackage[utf8]{inputenc}
\usepackage{amsmath,amssymb}
%\usepackage{amsthm} %\theoremstyle 
\usepackage{stmaryrd} %\ll/rrbracket
\usepackage{wasysym}
%\usepackage{dsfont}


%
\usepackage{etoolbox}

%
% Theorem environments; needs to be loaded before hyperref so that
% \autoref works
%
%\usepackage{amsthm,thmtools}
%
%\newcommand{\ownthmSpaceAbove}{5pt}
%\newcommand{\ownthmSpaceBelow}{5pt}
%% all thmtools hooks are within the theorem group, so we need to adjust global definitions
%\newcommand{\resetCurThmBraces}{%
%\gdef\curThmBraceOpen{(}%
%\gdef\curThmBraceClose{)}}
%\resetCurThmBraces
%\newcommand{\removeThmBraces}{%
%\gdef\curThmBraceOpen{}%
%\gdef\curThmBraceClose{}}
%
%\declaretheoremstyle[
%    spaceabove=\ownthmSpaceAbove,
%    spacebelow=\ownthmSpaceBelow,
%    headpunct=.,
%    postheadspace=.5em,
%    notebraces={\curThmBraceOpen}{\curThmBraceClose},
%    postheadhook={\resetCurThmBraces},
%]{definition}
%\declaretheoremstyle[
%    style=definition,
%    bodyfont=\itshape,
%    notebraces={\curThmBraceOpen}{\curThmBraceClose},
%    postheadhook={\resetCurThmBraces},
%]{theorem}

%
% To remove theorembrackets if needed

%

%\newcommand\mysubsecs[1]{\subparagraph*{#1}}

\newcommand{\resetCurThmBraces}{%
\gdef\curThmBraceOpen{(}%
\gdef\curThmBraceClose{)}}
\resetCurThmBraces
\newcommand{\removeThmBraces}{%
\gdef\curThmBraceOpen{}%
\gdef\curThmBraceClose{}}
\resetCurThmBraces

\newenvironment{notheorembrackets}{\removeThmBraces}{\resetCurThmBraces}

\usepackage{etoolbox}
\patchcmd{\thmhead}{(#3)}{\curThmBraceOpen #3\curThmBraceClose}{}{}


\usepackage{hyperref}
\hypersetup{hidelinks,final}

\usepackage[final,notref,notcite]{showkeys}
\usepackage[noadjust]{cite}


%%%% Text wrapping for showkeys
% Editorial wrapper omits seqsplit used only by inactive draft labels.
\usepackage{xstring}
\newcommand{\defaultshowkeysformat}[1]{%
% instead of \textvisiblespace you can also put in ~
% if you want to keep a plain space at space characters
\StrSubstitute{#1}{ }{\textvisiblespace}[\TEMP]%
\parbox[t]{\marginparwidth}{\raggedright\normalfont\small\ttfamily\(\{\){\color{red!50!black}\expandafter\seqsplit\expandafter{\TEMP}}\(\}\)}%
}
\newenvironment{hideshowkeys}{%
  \renewcommand*\showkeyslabelformat[1]{%
  \noexpandarg%
  }
}{%
  \renewcommand*\showkeyslabelformat[1]{%
  \noexpandarg%
  \defaultshowkeysformat{##1}%
  }
}

\renewcommand*\showkeyslabelformat[1]{%
\noexpandarg%
\defaultshowkeysformat{#1}%
}

\renewcommand\itemautorefname{Item}
\newcommand{\itemref}[2]{\autoref{#1}\ref{#2}}

% Symbols for bars
\newcommand{\longmid}{\hspace{0.9ex}\smash{\rule[-1.0ex]{0.41pt}{3.2ex}}\hspace{0.9ex}}
\newcommand{\midmid}{\hspace{0.2ex}{\rule[-0.1ex]{0.6pt}{1.65ex}}\hspace{0.2ex}}
\newcommand{\scriptmidmid}{\hspace{0.2ex}{\rule[-0.1ex]{0.6pt}{1.1ex}}\hspace{0.2ex}}
\newcommand{\subscriptmidmid}{\hspace{0.2ex}{\rule[-0.1ex]{0.6pt}{0.9ex}}\hspace{0.2ex}}

%Separator for set comprehension
\newcommand{\compr}{:}

\newcommand{\newletter}[1]{{\midmid}#1}
\newcommand{\scriptnew}[1]{{\scriptmidmid}#1}
\newcommand{\subscriptnew}[1]{{\subscriptmidmid}#1}

\newcommand{\trans}[1]{\xrightarrow{#1}}
\newcommand{\F}{\mathcal{F}}
\newcommand{\FN}{\mathsf{FN}}
\newcommand{\id}{\mathit{id}}

\renewcommand{\o}{\cdot}


\newcommand{\prelang}{L_\mathsf{pre}}
\newcommand{\nd}[1]{\mathsf{nd}(#1)}
\newcommand{\Names}[1]{\mathsf{N}(#1)}
\newcommand{\free}[1]{\mathsf{FN}(#1)}
\newcommand{\formulae}{\mathsf{Bar}}
\newcommand{\negepsilon}{\neg\epsilon}
\newcommand{\degree}{\mathsf{deg}}

\newcommand{\barstrings}{\barA^*}
\newcommand{\barA}{\overline{\names}}
\newcommand{\barname}{\beta}
\newcommand{\ub}{\mathsf{ub}}


%\usepackage{enumerate}
%
% Nice enumerate environments
%
% we shouldn't change the LIPIcs style; the paper is actually shorter with the
% standard envs
%
%\usepackage[inline]{enumitem}
%\setlist[enumerate,1]{label=(\arabic*),font=\normalfont,align=left,leftmargin=0pt,labelindent=0pt,listparindent=\parindent,labelwidth=0pt,itemindent=!,topsep=3pt,parsep=0pt,itemsep=3pt,start=1}
%\setlist[enumerate,2]{label=(\alph*),font=\normalfont,labelindent=*,leftmargin=*,start=1}
%\setlist[itemize]{labelindent=*,leftmargin=*,topsep=5pt,itemsep=3pt}
%\setlist[description]{labelindent=*,leftmargin=*,itemindent=-1 em}

\numberwithin{equation}{section}

% Editorial wrapper omits unused draft author-note package.
% Editorial draft author registration omitted: \FXRegisterAuthor{dh}{adh}{DH}
% Editorial draft author registration omitted: \FXRegisterAuthor{ls}{als}{LS}
% Editorial draft author registration omitted: \FXRegisterAuthor{sm}{asm}{SM}
% Editorial draft author registration omitted: \FXRegisterAuthor{hu}{ahu}{HU}

\usepackage{algorithmicx}
    \usepackage[ruled]{algorithm}
\usepackage[noend]{algpseudocode}


\usepackage{tikz}
 \usetikzlibrary{cd}
 \usetikzlibrary{trees}
 \usetikzlibrary{shapes}
 \usetikzlibrary{fit}
 \usetikzlibrary{shadows}
 \usetikzlibrary{backgrounds}
 \usetikzlibrary{arrows,automata,positioning}

\tikzset{
   ->,
   n/.style= {circle,fill,inner sep=1.5pt,node distance=2cm}
  ,acc/.style={circle,draw,inner sep=3pt,node distance=2cm}
  ,phantom/.style={circle},
  ,arr/.style={->, >=stealth, semithick, shorten <= 3pt, shorten >= 3pt}
}


\newcommand{\lesslex}{{<_\mathit{lex}}}

%
% make autoref work for sections
%
%\addto\extrasbritish{% only necessary if babel is used
  \renewcommand{\sectionautorefname}{Section}%
  \renewcommand{\subsectionautorefname}{Section}%
  %\renewcommand{\sectionname}{Section}%
%}


%\newcommand{\lesslex}{{<_\text{lex}}}

\newcommand{\CO}{\mathcal{O}}

\newcommand{\ol}{\overline}
\newcommand{\epito}{\twoheadrightarrow}
\newcommand{\monoto}{\hookrightarrow}
\newcommand{\xto}[1]{\xrightarrow{~#1~}}
\newcommand{\seq}{\subseteq}
\newcommand{\tl}{\widetilde}
\newcommand{\dash}{\mathop{-}}
\newcommand{\Perm}{\mathsf{Perm}}

\newcommand{\muBar}{\textsf{Bar-$\mu$TL}}
\newcommand{\states}{\mathsf{states}}
\newcommand{\prestates}{\mathsf{prestates}}
\newcommand{\quasistates}{\mathsf{quasistates}}
\renewcommand{\Box}{\square}
\renewcommand{\Diamond}{\lozenge}
\newcommand{\pow}{\mathcal{P}}
\newcommand{\hearts}{\heartsuit}
\newcommand{\prems}{\mathit{prems}}
\newcommand{\eval}{\mathit{eval}}
\newcommand{\uto}{\mathsf{uto}}
\newcommand{\mto}{\mathsf{mto}}
\newcommand{\maxdefs}{\mathit{maxbases}}
\newcommand{\boxfree}[1]{\mathsf{bf}(#1)}
\newcommand{\semb}[1]{(#1)}
\newcommand{\semS}[1]{\llbracket S\vdash #1 \rrbracket}
\newcommand{\sem}[1]{\llbracket #1 \rrbracket}
\newcommand{\litsem}[1]{\llparenthesis #1 \rrparenthesis}
\newcommand{\ind}[1]{\llparenthesis #1 \rrparenthesis}

\newcommand{\ad}{\mathsf{ad}}
\newcommand{\al}{\mathsf{al}}

\newcommand{\ie}{i.e.}

\newcommand{\To}{\Rightarrow}
\newcommand{\calB}{\mathcal{B}}
\newcommand{\calC}{\mathcal{C}}
\newcommand{\calU}{\mathcal{U}}
\newcommand{\calO}{\mathcal{O}}
\newcommand{\calW}{\mathcal{W}}
\newcommand{\calV}{\mathcal{V}}
\newcommand{\calI}{\mathcal{I}}
\newcommand{\calJ}{\mathcal{J}}
%\newcommand{\calN}{\mathcal{N}}
\newcommand{\calP}{\mathcal{P}}
\newcommand{\calF}{\mathcal{F}}

\newcommand{\sprop}{P}
\newcommand{\snom}{N}
\newcommand{\svar}{\mathsf{Var}}
\newcommand{\smod}{\Lambda}

\newcommand{\ALC}{\mathcal{ALC}}
\makeatletter
\def\moverlay{\mathpalette\mov@rlay}
\def\mov@rlay#1#2{\leavevmode\vtop{%
   \baselineskip\z@skip \lineskiplimit-\maxdimen
   \ialign{\hfil$\m@th#1##$\hfil\cr#2\crcr}}}
\newcommand{\charfusion}[3][\mathord]{
    #1{\ifx#1\mathop\vphantom{#2}\fi
        \mathpalette\mov@rlay{#2\cr#3}
      }
    \ifx#1\mathop\expandafter\displaylimits\fi}
\makeatother

\newcommand{\dotcup}{\charfusion[\mathbin]{\cup}{\cdot}}

\newcommand{\dom}{\mathsf{dom}}

\newcommand{\braket}[1]{\langle #1 \rangle}

\renewcommand{\epsilon}{\varepsilon}

\newcommand{\PSPACE}{\mathrm{PSPACE}}
\newcommand{\orb}{\mathsf{orb}}
\newcommand{\fs}{\mathsf{fs}}
\newcommand\pfun{\mathrel{\ooalign{\hfil$\mapstochar\mkern5mu$\hfil\cr$\to$\cr}}}
\newcommand{\rank}{\mathsf{rk}}
\newcommand{\argument}{\_\!\!\_}
\newcommand{\Rules}{\mathcal R}
\newcommand{\fpvars}{\mathbf{V}}
\newcommand{\concl}{\mathit{Cn}}
\newcommand{\eloise}{\mathsf{Eloise}}
\newcommand{\abelard}{\mathsf{Abelard}}
\newcommand{\track}{\mathsf{track}}
\newcommand{\detcarrier}{D_\target}
\newcommand{\detprio}{\beta}
\newcommand{\FL}{\mathsf{FL}}
\newcommand{\size}{\mathsf{size}}
\newcommand{\Rec}{\mathit{Inh}}
\newcommand{\defer}{\mathsf{dfr}}
\newcommand{\PLentails}{\vdash_{\mathit{PL}}}
\newcommand{\target}{\chi}
\newcommand{\FLtarget}{\mathbf F}
\newcommand{\FLphi}{\mathbf F}
\newcommand{\snodes}{\mathbf{S}}
\newcommand{\nodes}{\mathbf{N}}
\newcommand{\fnodes}{\mathbf{C}}
\newcommand{\deferrals}{\mathbf{D}}
\newcommand{\Cls}{\mathcal}
\newcommand{\CM}{{\Cls M}}
\newcommand{\FA}{{\mathfrak A}}
\newcommand{\ATLdiamond}[1]{{\langle\!\langle#1\rangle\!\rangle}}
\newcommand{\ATLbox}[1]{{[\![#1]\!]}}
\newcommand{\until}{U}
\newcommand{\Dist}{\mathcal{D}}
\newcommand{\Rat}{{\mathbb{Q}}}
\newcommand{\Nat}{{\mathbb{N}}}
\newcommand{\Real}{{\mathbb{R}}}
\newcommand{\Bag}{\mathcal{B}}
\newcommand{\Gm}{\mathcal{G}}
\newcommand{\gldiamond}[1]{\Diamond_{#1}}
\newcommand{\glbox}[1]{\square_{#1}}
\newcommand{\lsem}{\llbracket}
\newcommand{\rsem}{\rrbracket}
\newcommand{\bbrack}[1]{\llbracket #1 \rrbracket}
\newcommand{\cset}[1]{{\{ #1\} }}
\newcommand{\tup}[1]{{\langle #1 \rangle}}
\newcommand{\pa}[1]{{\left( #1 \right)}}
\newcommand{\nat}{\mathbb{N}}
\newcommand\pto{\mathrel{\ooalign{\hfil$\mapstochar$\hfil\cr$\to$\cr}}}
\newcommand\NExpTime{$\textsc{NExpTime}$\xspace}
\newcommand\NExpSpace{$\textsc{NExpSpace}$\xspace}
\newcommand{\dual}[1]{\overline{#1}}
\newcommand{\FLang}{\mathcal{F}}	
\newcommand{\otto}{\leftrightarrow}	
\newcommand{\Set}{\mathsf{Set}}
\newcommand{\Ni}{\noindent}
\newcommand{\Sem}[1]{{[\![#1]\!]}}
\newcommand{\fe}{{\forall\exists}}
\newcommand{\ef}{{\exists\forall}}
\newcommand{\ha}{{\hearts\alpha}}
\newcommand{\psem}[1]{\widehat{[\![#1]\!]}}
\newcommand{\contrapow}{\breve{P}}
\newcommand{\Op}{\mathit{op}}
\newcommand{\lO}{\mathcal{O}}
\newcommand{\infrule}[2]{\frac{#1}{#2}}
\newcommand{\limp}{\Rightarrow}
\newcommand{\+}[1]{\mathcal{#1}\xspace}
\newcommand{\sut}{such that\xspace}
\newcommand{\ov}[1]{\overline{#1}}

\newcommand{\op}{{\mathsf{op}}}
\newcommand{\Pow}{\mathcal{P}}
\newcommand{\ufs}{{\mathsf{ufs}}}

\newcommand{\At}{\mathsf{At}}
\newcommand{\Act}{\mathsf{Act}}
\DeclareMathOperator{\Nom}{\mathbf{Nom}}
\DeclareMathOperator{\LFP}{\mathsf{LFP}}
\DeclareMathOperator{\FP}{\mathsf{FP}}
\DeclareMathOperator{\GFP}{\mathsf{GFP}}
\DeclareMathOperator{\inv}{\mathsf{inv}}
\DeclareMathOperator{\univ}{univ}
\DeclareMathOperator{\ga}{ga}
\DeclareMathOperator{\cl}{cl}
\DeclareMathOperator{\flc}{\mathsf{FL}}
\DeclareMathOperator{\chr}{char}
\DeclareMathOperator{\classicalEF}{\mathsf{EF}}
\DeclareMathOperator{\classicalAG}{\mathsf{AG}}
\DeclareMathOperator{\tout}{to}
\newcommand{\reach}{{R_\infty}}

\newcommand{\N}{\mathbb{N}}
\newcommand{\Q}{\mathbb{Q}}

%bold E and A in mathmode
\newcommand{\E}{\mathbf{E}}
\newcommand{\A}{\mathbf{A}}

\newcommand{\EF}{\classicalEF_c}
\newcommand{\AG}{\classicalAG_c}

%equivalence diamond and box
\newcommand{\equivdia}{\langle \equiv \rangle}
\newcommand{\equivbox}{[ \equiv ]}

\newcommand{\restrict}{\vert}

%
% From Stefan
%

\newcommand{\takeout}[1]{\empty}
\newcommand{\names}{\mathbb{A}}
\DeclareMathOperator{\fix}{\mathsf{fix}}
\DeclareMathOperator{\Fix}{\mathsf{Fix}}
\DeclareMathOperator{\supp}{\mathsf{supp}} %% spacing mit Betragsstrichen komisch
%\newcommand{\supp}{\mathsf{supp}}
\DeclareMathOperator{\clos}{\mathsf{cl}}
\newcommand{\fresh}{\mathbin{\#}}
\newcommand{\powufs}{\pow_{\ufs}}
\newcommand{\powfs}{\pow_{\fs}}
\newcommand{\powf}{\pow_{\omega}}


% for Springer LLNCS
%\makeatletter
%\newcommand\mysubsec{\@startsection{paragraph}{4}{\z@}%
%  {-6\p@ \@plus -4\p@ \@minus -4\p@}%
%  {-0.5em \@plus -0.22em \@minus -0.1em}%
%  {\normalfont\normalsize\bfseries}}
%\makeatother

% For LIPIcs
%\makeatletter
%\newcommand{\mysubsec}[1]{%
%  \par\medskip\noindent{\bfseries #1}\hspace{3mm}%
%  \@ifnextchar\par{\@gobble}{}% this eats a following \par if any
%}
%\makeatother
\newcommand{\mysubsec}{\subparagraph*}

%
% Extra theorem envs
%

% LIPIcs stuff
%\theoremstyle{plain}


\theoremstyle{definition}
\newtheorem{defn}[theorem]{Definition}
\newtheorem{expl}[theorem]{Example}
\newtheorem{rem}[theorem]{Remark}
\newtheorem{nota}[theorem]{Notation}
\newtheorem{ass}[theorem]{Assumption}
\newtheorem*{fact}{Fact}



%
% Theorem envs
%
%\theoremstyle{theorem}
%\newtheorem{theorem}{Theorem}[section]
%\newtheorem{corollary}[theorem]{Corollary}
%\newtheorem{lemma}[theorem]{Lemma}
%\newtheorem{proposition}[theorem]{Proposition}

%\theoremstyle{definition}
%\newtheorem{defn}[theorem]{Definition}
%\newtheorem{expl}[theorem]{Example}
%\newtheorem{rem}[theorem]{Remark}
%\newtheorem{nota}[theorem]{Notation}
%\newtheorem{ass}[theorem]{Assumption}
%\newtheorem{fact}{Fact}


\newcommand\ExpTime{$\textsc{ExpTime}$\xspace}
\newcommand\ExpSpace{$\textsc{ExpSpace}$\xspace}
\newcommand\PSpace{$\textsc{PSpace}$\xspace}
\newcommand\NP{$\textsc{NP}$\xspace}
\newcommand\coNP{$\textsc{co-NP}$\xspace}

%
% Frontmatter
%
